% --- ARCHIVO: Capitulos/combate.tex ---

\chapter{El Arte de la Guerra}

\begin{loreblock}
``No es el arma, es la mano que la sostiene. No es la mano, es la voluntad detrás.''
\end{loreblock}

\vspace{0.4cm}

\section{El Pool de Ataque}

Cuando una unidad declara un ataque \textbf{Melee} o \textbf{Rango}, se lanza un \textit{pool} de tres dados simultáneamente. Hacerlo al mismo tiempo acelera el juego.

\begin{reglaespecial}{La Tirada de Triple Dado}
\begin{itemize}
    \item \textbf{1d12 (Éxito):} Determina si el ataque impacta. Se suma al atributo atacante y se enfrenta a la defensa del objetivo.
    \item \textbf{1d100 (Localización):} Determina en qué parte del cuerpo impacta el golpe, \textit{solo si el ataque tuvo éxito}.
    \item \textbf{1d6 (Efecto/Estado):} Si el ataque aplica un estado negativo (sangrado, veneno, etc.), este dado determina la potencia del estado.
\end{itemize}
\end{reglaespecial}

\section{Atributos Principales (Escala 1--6)}

Todos los atributos tienen un valor entre +1 y +6 y se suman a la tirada de 1d12. Los modificadores de equipo y situación pueden añadir de -4 a +4 adicionales. \textbf{Ningún valor puede ser menor a 0}.

\begin{center}
\begin{tabularx}{\linewidth}{|l|X|}
\hline
\rowcolor{HammerRed!15} \textbf{Atributo} & \textbf{Uso} \\ \hline
\textbf{Melee ATK} & Precisión en combate cuerpo a cuerpo. \\ \hline
\textbf{Range ATK} & Precisión con armas a distancia. \\ \hline
\textbf{Reflejos} & Esquiva. Dificultad para ser impactado. Cuesta 2 PA como Reacción. \\ \hline
\textbf{Bloqueo} & Defensa. Dificultad para ser impactado con equipo/técnica. \\ \hline
\textbf{Focus} & Mantener o potenciar efectos mágicos/técnicos. \\ \hline
\textbf{Res. Física} & Bono a 1d6 contra estados físicos (sangrado, veneno). \\ \hline
\textbf{Res. Mágica} & Bono a 1d6 contra estados mágicos (sueño, miedo). \\ \hline
\end{tabularx}
\end{center}

\section{Resolución Paso a Paso}

\begin{reglaespecial}{Paso 1: Tirada de Éxito (Ataque vs. Defensa)}
El atacante tira: \textbf{1d12 + Atributo de Ataque + Modificadores}

El defensor responde con una tirada de \textbf{Salvación}:
\begin{itemize}
    \item Contra Melee: \textbf{1d12 + Bloqueo + Modificadores}
    \item Contra Rango: \textbf{1d12 + Reflejos + Modificadores} (cuesta 2 PA — ver Reacciones)
\end{itemize}

Si el resultado del \textbf{atacante} es \textbf{igual o mayor} al del defensor: \textbf{el ataque impacta}.
\end{reglaespecial}

\begin{reglaespecial}{Paso 2: Localización (1d100)}
Si impacta, consulta la tabla. El resultado del dado ya fue tirado junto al pool inicial.
\end{reglaespecial}

\begin{reglaespecial}{Paso 3: Resolución del Estado Negativo (si aplica)}
Si el ataque puede infligir un estado (sangrado, veneno, sueño, etc.):

\textbf{Atacante:} 1d6 + Modificadores de Arma\\
\textbf{Defensor:} 1d6 + Resistencia correspondiente (Física o Mágica)

Si el atacante supera al defensor: \textbf{el estado se aplica}.
\end{reglaespecial}

\subsection*{Ejemplo Completo de Combate}

\begin{ejemplo}{Ejemplo: Ataque Melee con posibilidad de Sangrado}
\textbf{Jugador A} ataca con su unidad (Melee ATK +3, con modificador +1 por flanco). El objetivo del \textbf{Jugador B} tiene Bloqueo +2 y Modificador de Escudo +1.

Jugador A tira simultáneamente:
\begin{itemize}
    \item \textbf{Éxito:} $1d12 + 3 + 1 = [6] + 4 = 10$
    \item \textbf{Localización:} $1d100 = 48$
    \item \textbf{Efecto (Sangrado):} $1d6 + 1 = [4] + 1 = 5$
\end{itemize}
Jugador B responde:
\begin{itemize}
    \item \textbf{Salvación:} $1d12 + 2 + 1 = [4] + 3 = 7$
    \item \textbf{Resistencia Física:} $1d6 + 2 = [3] + 2 = 5$
\end{itemize}
\textbf{Resultado:}
\begin{enumerate}
    \item Ataque (10) $>$ Defensa (7): \textbf{Impacta}.
    \item Localización 48 → \textbf{Cuerpo}.
    \item Estado: Efecto (5) \textbf{no supera} Resistencia (5) → El sangrado no se aplica.
    \item Daño: El arma hace 4 de daño. Cuerpo tiene RD 2. \textbf{2 PV de daño al Cuerpo}.
\end{enumerate}
\end{ejemplo}

\section{Localización de Daño}

\begin{imagebox}{Diagrama de un personaje con las zonas de localización señaladas — Cabeza, Cuerpo, Brazos, Piernas}
    \vspace{3.5cm}
\end{imagebox}

\begin{tabularx}{\linewidth}{|l|c|c|c|X|}
\hline
\rowcolor{HammerRed!20} \textbf{Zona} & \textbf{1d100} & \textbf{PV Máx.} & \textbf{RD Base} & \textbf{Efecto al llegar a 0 PV} \\ \hline
\textbf{Cabeza} & 96--100 & 4 & 0/1/2/3/4 & \textbf{Incapacitado} o Muerte (según facción) \\ \hline
\textbf{Brazos} & 81--95 & 6 & 0/1/2/3/4 & Incapaz de usar armas en ese brazo \\ \hline
\textbf{Piernas} & 66--80 & 6 & 0/1/2/3/4 & Movimiento reducido al 50\% \\ \hline
\textbf{Cuerpo} & 01--65 & 10 & 0/1/2/3/4 & \textbf{UNIDAD ELIMINADA} \\ \hline
\end{tabularx}

\vspace{0.3cm}

\subsection*{Reducción de Daño (RD)}

La RD representa la resistencia de la armadura y la toughness natural de la unidad en esa zona. El daño final recibido es:

\begin{tcolorbox}[colback=HammerRed!8, colframe=HammerRed]
\centering
\Large \textbf{Daño Final = Daño Entrante $-$ RD de la Zona}
\end{tcolorbox}

\vspace{0.2cm}

La RD de cada zona está determinada por las \textbf{características de la unidad y su equipamiento}. Los modificadores van de -4 a +4. \textbf{Ningún valor puede bajar de 0.}

\begin{reglaespecial}{Incapacitación}
Cuando una zona llega a 0 PV, su efecto se aplica de forma \textbf{permanente} por el resto de la partida:
\begin{itemize}
    \item \textbf{Cabeza:} La unidad queda Incapacitada (pierde toda acción) o muere directamente, según las reglas de su facción.
    \item \textbf{Brazos:} No puede usar el arma equipada en ese brazo. Si ambos brazos están a 0, no puede atacar en Melee ni Rango.
    \item \textbf{Piernas:} MOV reducido permanentemente a la mitad. Redondeado hacia arriba.
    \item \textbf{Cuerpo:} La unidad es eliminada y cae en el tablero.
\end{itemize}
\end{reglaespecial}

\section{Estados Negativos}

Los estados negativos son efectos que persisten entre rondas o durante un tiempo determinado. Se indica en cada habilidad cuándo expiran.

\begin{center}
\begin{tabularx}{\linewidth}{|l|X|}
\hline
\rowcolor{GoldRule!20} \textbf{Estado} & \textbf{Efecto} \\ \hline
\textbf{Sangrado} & Al inicio de cada ronda, pierde 1 PV de la zona afectada. \\ \hline
\textbf{Envenenado} & Al inicio de cada ronda, pierde 1 PV de Cuerpo. Se apila hasta x4. \\ \hline
\textbf{Derribado} & Gasta 1 PA extra para levantarse. Melee enemigo tiene +2 mientras esté en el suelo. \\ \hline
\textbf{Enmarañado} & No puede moverse. $-1$ a CA, Reflejos y Bloqueo. Puede intentar liberarse con Resistencia Física. \\ \hline
\textbf{Aturdido} & Pierde 1 PA en su próxima activación. \\ \hline
\textbf{Dormido} & No puede activarse. Se despierta si recibe daño o al inicio de su turno con una tirada de Resistencia Mágica. \\ \hline
\end{tabularx}
\end{center}

\section{Ataques de Distancia y Pérdida de Precisión}

Los ataques de rango tienen penalizadores acumulativos según la distancia:

\begin{itemize}
    \item \textbf{Rango corto} (0--50\% del rango del arma): Sin penalizador.
    \item \textbf{Rango medio} (50--75\%): $-1$ a la tirada de Éxito.
    \item \textbf{Rango largo} (75--100\%): $-2$ a la tirada de Éxito.
    \item \textbf{Fuera de rango:} No se puede disparar.
\end{itemize}

Cada arma puede tener sus propias penalizaciones específicas indicadas en su ficha.
